
\documentclass{article}
\usepackage{colortbl}
\usepackage{makecell}
\usepackage{multirow}
\usepackage{supertabular}

\begin{document}

\newcounter{utterance}

\centering \large Interaction Transcript for game `wordle', experiment `high\_frequency\_words\_no\_clue\_no\_critic', episode 1 with mistral{-}large{-}3{-}azure.
\vspace{24pt}

{ \footnotesize  \setcounter{utterance}{1}
\setlength{\tabcolsep}{0pt}
\begin{supertabular}{c@{$\;$}|p{.15\linewidth}@{}p{.15\linewidth}p{.15\linewidth}p{.15\linewidth}p{.15\linewidth}p{.15\linewidth}}
    \# & \multicolumn{2}{c}{Player} && \multicolumn{2}{c}{Game Master} \\
    \hline

    \theutterance \stepcounter{utterance}  
    & & & \multicolumn{4}{p{0.6\linewidth}}{
        \cellcolor[rgb]{0.9,0.9,0.9}{
            \makecell[{{p{\linewidth}}}]{
                \texttt{\tiny{[P1$\langle$GM]}}
                \texttt{Ви — мовний чарівник, який любить вгадувати слова, використовуючи задані правила.} \\
\\ 
\texttt{Ласкаво просимо до Wordle! У вас є шість спроб, щоб вгадати цільове слово — допустиме слово з п'яти малих літер. Будь ласка, використовуйте теги "пояснення:" та "спроба:", щоб надати стислий опис для кожної спроби.} \\
\\ 
\texttt{Наприклад, якщо ваша спроба — "книга", ваша відповідь має бути} \\
\texttt{пояснення: це поширене слово з п'яти літер, і я починаю зі спроби з цього слова.} \\
\texttt{спроба: книга} \\
\\ 
\texttt{Після кожної спроби вашу відповідь буде перевірено, і ви отримаєте відгук, який вказує, які літери правильні (зелений), які літери правильні, але стоять не на своєму місці (жовтий), і які літери неправильні (червоний). Цей відгук може бути корисним для визначення, які літери слід включити або виключити у вашій наступній спробі.} \\
\\ 
\texttt{Наприклад, відгук для "книга" може бути:} \\
\texttt{відгук\_на\_спробу: к<жовтий> н<жовтий> и<зелений> г<жовтий> а<червоний>} \\
\\ 
\texttt{Пояснення має містити деталі про те, як відгук\_на\_спробу використовується, щоб дійти до нової спроби.} \\
\\ 
\texttt{Не генеруйте жодного іншого тексту.} \\
\\ 
\texttt{Почнімо з вашої першої спроби.} \\
            }
        }
    }
     \\ \\

    \theutterance \stepcounter{utterance}  
    & \multicolumn{4}{p{0.6\linewidth}}{
        \cellcolor[rgb]{0.561,0.737,0.561}{
            \makecell[{{p{\linewidth}}}]{
                \texttt{\tiny{[P1$\rangle$GM]}}
                \texttt{пояснення: Виберу поширене слово з п'яти літер, яке містить різноманітні голосні та приголосні, щоб максимально звузити можливості. Слово "апель" підходить, оскільки включає поширені літери.} \\
\texttt{спроба: апель} \\
            }
        }
    }
    & & \\ \\

    \theutterance \stepcounter{utterance}  
    & & & \multicolumn{2}{p{0.3\linewidth}}{
        \cellcolor[rgb]{0.9,0.9,0.9}{
            \makecell[{{p{\linewidth}}}]{
                \texttt{\tiny{[GM$|$GM]}}
                \texttt{attempts = 1} \\
\texttt{target\_word = право} \\
\texttt{guess = апель} \\
\texttt{guess\_feedback = а<yellow> п<yellow> е<red> л<red> ь<red>} \\
            }
        }
    }
    & & \\ \\

    \theutterance \stepcounter{utterance}  
    & & & \multicolumn{4}{p{0.6\linewidth}}{
        \cellcolor[rgb]{0.9,0.9,0.9}{
            \makecell[{{p{\linewidth}}}]{
                \texttt{\tiny{[P1$\langle$GM]}}
                \texttt{відгук\_на\_спробу: а<жовтий> п<жовтий> е<червоний> л<червоний> ь<червоний>} \\
\\ 
\texttt{Надайте відповідь лише в цьому форматі.} \\
\texttt{пояснення: деталі} \\
\texttt{спроба: слово} \\
            }
        }
    }
     \\ \\

    \theutterance \stepcounter{utterance}  
    & \multicolumn{4}{p{0.6\linewidth}}{
        \cellcolor[rgb]{0.561,0.737,0.561}{
            \makecell[{{p{\linewidth}}}]{
                \texttt{\tiny{[P1$\rangle$GM]}}
                \texttt{пояснення: Літера "а" є в слові, але не на першій позиції (жовтий), "п" також присутня, але не на другій (жовтий). Літери "е", "л", "ь" відсутні (червоний). Вибираю слово з "а" і "п" на інших позиціях, без "е", "л", "ь". Наприклад, "парус".} \\
\texttt{спроба: парус} \\
            }
        }
    }
    & & \\ \\

    \theutterance \stepcounter{utterance}  
    & & & \multicolumn{2}{p{0.3\linewidth}}{
        \cellcolor[rgb]{0.9,0.9,0.9}{
            \makecell[{{p{\linewidth}}}]{
                \texttt{\tiny{[GM$|$GM]}}
                \texttt{attempts = 2} \\
\texttt{target\_word = право} \\
\texttt{guess = парус} \\
\texttt{guess\_feedback = п<green> а<yellow> р<yellow> у<red> с<red>} \\
            }
        }
    }
    & & \\ \\

    \theutterance \stepcounter{utterance}  
    & & & \multicolumn{4}{p{0.6\linewidth}}{
        \cellcolor[rgb]{0.9,0.9,0.9}{
            \makecell[{{p{\linewidth}}}]{
                \texttt{\tiny{[P1$\langle$GM]}}
                \texttt{відгук\_на\_спробу: п<зелений> а<жовтий> р<жовтий> у<червоний> с<червоний>} \\
\\ 
\texttt{Надайте відповідь лише в цьому форматі.} \\
\texttt{пояснення: деталі} \\
\texttt{спроба: слово} \\
            }
        }
    }
     \\ \\

    \theutterance \stepcounter{utterance}  
    & \multicolumn{4}{p{0.6\linewidth}}{
        \cellcolor[rgb]{0.561,0.737,0.561}{
            \makecell[{{p{\linewidth}}}]{
                \texttt{\tiny{[P1$\rangle$GM]}}
                \texttt{пояснення: Літера "п" стоїть на правильній позиції (зелений), "а" є в слові, але не на другій (жовтий), "р" також присутня, але не на третій (жовтий). Літери "у" та "с" відсутні (червоний). Вибираю слово з "п" на першій позиції, "а" на іншій, "р" на іншій, без "у" та "с". Наприклад, "парта".} \\
\texttt{спроба: парта} \\
            }
        }
    }
    & & \\ \\

    \theutterance \stepcounter{utterance}  
    & & & \multicolumn{2}{p{0.3\linewidth}}{
        \cellcolor[rgb]{0.9,0.9,0.9}{
            \makecell[{{p{\linewidth}}}]{
                \texttt{\tiny{[GM$|$GM]}}
                \texttt{attempts = 3} \\
\texttt{target\_word = право} \\
\texttt{guess = парта} \\
\texttt{guess\_feedback = п<green> а<yellow> р<yellow> т<red> а<red>} \\
            }
        }
    }
    & & \\ \\

    \theutterance \stepcounter{utterance}  
    & & & \multicolumn{4}{p{0.6\linewidth}}{
        \cellcolor[rgb]{0.9,0.9,0.9}{
            \makecell[{{p{\linewidth}}}]{
                \texttt{\tiny{[P1$\langle$GM]}}
                \texttt{відгук\_на\_спробу: п<зелений> а<жовтий> р<жовтий> т<червоний> а<червоний>} \\
\\ 
\texttt{Надайте відповідь лише в цьому форматі.} \\
\texttt{пояснення: деталі} \\
\texttt{спроба: слово} \\
            }
        }
    }
     \\ \\

    \theutterance \stepcounter{utterance}  
    & \multicolumn{4}{p{0.6\linewidth}}{
        \cellcolor[rgb]{0.561,0.737,0.561}{
            \makecell[{{p{\linewidth}}}]{
                \texttt{\tiny{[P1$\rangle$GM]}}
                \texttt{пояснення: "П" на першій позиції правильно (зелений). "А" є в слові, але не на другій (жовтий), "р" також присутня, але не на третій (жовтий). Літери "т" та друга "а" відсутні (червоний). Отже, "а" має бути лише один раз, а "р" не на 3{-}й позиції. Вибираю слово з "п" на 1{-}й, "а" на 4{-}й або 5{-}й, "р" на 2{-}й або 4{-}й. Наприклад, "право".} \\
\texttt{спроба: право} \\
            }
        }
    }
    & & \\ \\

    \theutterance \stepcounter{utterance}  
    & & & \multicolumn{2}{p{0.3\linewidth}}{
        \cellcolor[rgb]{0.9,0.9,0.9}{
            \makecell[{{p{\linewidth}}}]{
                \texttt{\tiny{[GM$|$GM]}}
                \texttt{attempts = 4} \\
\texttt{target\_word = право} \\
\texttt{guess = право} \\
\texttt{guess\_feedback = п<green> р<green> а<green> в<green> о<green>} \\
            }
        }
    }
    & & \\ \\

    \theutterance \stepcounter{utterance}  
    & & & \multicolumn{2}{p{0.3\linewidth}}{
        \cellcolor[rgb]{0.9,0.9,0.9}{
            \makecell[{{p{\linewidth}}}]{
                \texttt{\tiny{[GM$|$GM]}}
                \texttt{game\_result = WIN} \\
            }
        }
    }
    & & \\ \\

\end{supertabular}
}

\end{document}
